\documentclass[10pt,a4paper]{amsart}
\usepackage[margin=19mm]{geometry}
\usepackage{amsmath,amssymb,mathtools}
\usepackage[scr=rsfs,cal=euler]{mathalfa}
\usepackage{xfrac}
\usepackage[unicode,hidelinks,bookmarks=false]{hyperref}
\newcommand{\ie}{i.e.\ }
\newcommand{\cf}{cf.\ }
\newcommand{\resp}{respectively\ }
\newcommand{\Subsection}{\subsection}
\providecommand{\nobreakdash}{\nobreak\hbox{-}}
\setlength{\emergencystretch}{2em}
\multlinegap0pt
\usepackage{amssymb}
\usepackage{dsfont}
\usepackage{algorithm}\usepackage{algpseudocode}
\newtheorem{lemma}{Lemma}
\newcommand{\mat}[1]{\mathbf{#1}}
\renewcommand{\vec}[1]{\mathbf{#1}}
\title{Accelerating a restarted Krylov method for matrix functions with randomization}
\author{Nicolas L. Guidotti, Per-Gunnar Martinsson, Juan A. Acebr\'on, Jos\'e Monteiro}
\date{Source: arXiv:2503.22631v5 (CC BY 4.0)}
\begin{document}
\maketitle

\noindent\emph{Source and licence.} Accelerating a restarted Krylov method for matrix functions with randomization. Version arXiv:2503.22631v5 (CC BY 4.0), distributed by arXiv under the Creative Commons Attribution 4.0 International licence, http://creativecommons.org/licenses/by/4.0/, as recorded in the arXiv licence metadata for this exact version and shown on the abstract page https://arxiv.org/abs/2503.22631v5. Full licence notice: \texttt{licenses/arxiv-2503.22631-CC-BY-4.0.txt}.\par\smallskip

\noindent\textit{Editorial scope.} Counted unit: the article's lemma theorem:similarity (source0.tex lines 208--215). The article's complete original proof of it is delivered with it (source0.tex lines 216--226, one \texttt{proof} environment of 11 lines). No mathematics is written, repaired or re-derived: the statement and the proof are the article's own lines, in the article's own notation, and no retained line is substituted or re-flowed.  Retained uncounted as the source-local context the proof needs: lines 142--145.  Nothing was omitted from any retained span, so the package carries no omission record.  No retained line contains a citation, so the wrapper carries no bibliography and no bibliography entry.  No retained line contains a figure environment, an image-inclusion command or drawing code, so no image is delivered and the package declares no asset dependency.  Editorial formatting only, in addition: an AMS \texttt{amsart} preamble that restates the article's own declarations and macros used by the retained lines, namely the environments \texttt{lemma} on separate counters, together with the standard packages those lines need.  The retained environments print in this document as Lemma 1 in order of appearance, whereas the article prints the same items with the numbering of its own full document.  The complete native source of the article as distributed in the arXiv e-print for this exact version is bundled unchanged as \texttt{sources/175570/source0.tex}.
\begin{equation}
\label{eq:rand_arnoldi}
\mat{A} \mat{W}_m = \mat{W}_{m + 1} \mat{\underline{R}}_m =  \mat{W}_{m} \mat{R}_m + \vec{w}_{m + 1} \, r_{m + 1, m} \, \vec{e}^T_m,
\end{equation}

\begin{lemma}
\label{theorem:similarity}
As both $\mat{V}_m$ and $\mat{W}_m$ are bases for the same subspace $\mathcal{K}_m(\mat{A}, \vec{b})$, there is a nonsingular, upper triangular matrix $\mat{Z} \in \mathbb{C}^{m \times m}$ such that $\alpha \mat{Z}\vec{e}_1 = \beta \vec{e}_1$ and $\mat{W}_m = \mat{V}_m \mat{Z}$. Then,
\begin{equation}
\label{eq:similarity}
\mat{Z} \mat{R}_m \mat{Z}^{-1} = \mat{H}_m - \frac{r_{m+1, m}}{z_{mm}} \, \vec{c} \, \vec{e}_m^T, \quad \vec{c} = \mat{V}_m^* \vec{w}_{m + 1}.
\end{equation}
\end{lemma}

\begin{proof}
   Replacing $\mat{W}_m$ with $\mat{V}_m \mat{Z}$ in the Arnoldi-like decomposition (\ref{eq:rand_arnoldi}), we have
   \begin{equation*}
     \mat{A} \mat{V}_m \mat{Z} = \mat{V}_m \mat{Z} \mat{R}_m + \vec{w}_{m + 1} \, r_{m + 1, m} \, \vec{e}^T_m.
   \end{equation*}
   Multiplying by $\mat{V}_m^*$ on the left and by $\mat{Z}^{-1}$ on the right,
      \begin{equation*}
     \mat{V}_m^* \mat{A} \mat{V}_m =  \mat{Z} \mat{R}_m \mat{Z}^{-1} + (\mat{V}_m^* \vec{w}_{m + 1}) \, r_{m + 1, m} \, \vec{e}^T_m \mat{Z}^{-1}.
   \end{equation*}
   By construction, $\mat{H}_m = \mat{V}_m^* \, \mat{A} \mat{V}_m$ and $\vec{e}_m^T \mat{Z}^{-1} = 1 / z_{mm} \vec{e}_m^T$ since $\mat{Z}$ is upper triangular. Setting $\vec{c} = \mat{V}_m^* \vec{w}_{m + 1}$ and rearranging the terms completes the proof.
\end{proof}

\end{document}
